\documentclass{article}
\usepackage{graphicx} % Required for inserting images
\usepackage[margin=1in]{geometry}
\usepackage{float} 
\usepackage{amsmath}
\usepackage{xcolor}



\begin{document}

\title{Fluid and Biogeochemical Dynamics in Sinking Marine Particle}
\author{Conor Padmanabhan}
\date{March 2026}


\maketitle

\tableofcontents

\section{Introduction}

There's plenty of oxygen in the ocean. However, within particles, they can run out of oxygen, and other processes begin. N$_2$O (GHG) contribution: Wan et al. (2023) paper is good for background of specifically this..


\section{Science Questions}
\begin{itemize}
    \item Effect of sinking speed on the amount of $N_2 O$ leakage?
    \item How does the size of the anoxic core change? Definitely from O$_2$ levels. But what about sinking speed, size..
    \item For MicrOMZ: In the wake, could microbial species be prevalent? Like hotspots in little swirls.
    \item Look at coastal nitrogen loss ($N_2 O \& N_2$) in different $O_2$ regimes?
    \item \textcolor{gray}{As respiration consumes O$_2$, even in oxygenated waters, can the inside of a particle become anoxic, causing increased activity of N$_2$O-producing microbes (that survive without O$_2$)?}
\end{itemize}


\section{Model Formulation}
The physical system is modeled using the 2D incompressible Navier-Stokes equations, transformed into the vorticity-streamfunction formulation to enforce incompressibility implicitly and eliminate the pressure gradient term.

When O$_2$ is available, implicit microbes (aerobic heterotrophs) within the particle perform respiration, consuming O$_2$. When O$_2$ drops below a certain threshold, anaerobic heterotrophs start to take over, consuming nitrogen within the particle and eventually producing N$_2$O.

\subsection{Governing Fluid Equations}
To model a stationary, porous particle without relying on complex boundary-fitted grids or internal no-slip walls, an immersed boundary approach is utilized. A localized linear drag force, $\vec{F} = -\alpha(x,y) \vec{v}$, is applied, where $\alpha(x,y)$ represents the spatially varying drag coefficient that is strictly positive within the particle and zero in the free stream.

Taking the curl of the momentum equations brings in the vorticity. These are the equations\textsuperscript{1}:
\begin{equation}
    \frac{\partial \omega_z}{\partial t} - J(\psi, \omega_z) = \nu \nabla^2 \omega_z - \alpha \omega_z + u \frac{\partial \alpha}{\partial y} - v \frac{\partial \alpha}{\partial x}
\end{equation}
where $\nu=0.000001$ is the kinematic viscosity, $\omega_z$ is the vertical component of vorticity, and $J(\psi,\omega_z)$ is the Arakawa Jacobian:
\begin{align*}
    J_{a}(a, b) &= \frac{1}{3} \left[ J(a, b) + \hat{J}(a, b) + \tilde{J}(a, b) \right] \\
\hat{J}(a, b) &= \frac{\partial}{\partial x} \left( a \frac{\partial b}{\partial y} \right) - \frac{\partial}{\partial y} \left( a \frac{\partial b}{\partial x} \right) \\
\tilde{J}(a, b) &= \frac{\partial}{\partial y} \left( b \frac{\partial a}{\partial x} \right) - \frac{\partial}{\partial x} \left( b \frac{\partial a}{\partial y} \right)
\end{align*}


The velocity components $(u,v)$ are recovered by solving the Poisson equation for the streamfunction $\psi$:
\begin{equation}
    \nabla^2 \psi = -\omega_z
\end{equation}
where
\begin{equation}
    (u,v) = \left(-\frac{\partial \psi}{\partial y}, +\frac{\partial \psi}{\partial x}\right)
\end{equation}
To simulate an open channel with an incoming current, the streamfunction is decomposed into a uniform background flow and a perturbation: $\psi_{\text{total}} = U y + \psi'$.



\subsection{Passive Tracer Transport}
The concentration of any tracer, $C$, is affected by advection by the fluid and diffusion at a constant rate $K$. The transport equation is given by:
\begin{equation}
    \frac{\partial C}{\partial t} =- \vec{v} \cdot \nabla C = K \nabla^2 C
\end{equation}
The numerical scheme uses an enstrophy-conserving Arakawa Jacobian for the advective terms and a standard 5-point central difference stencil for the Laplacian diffusion terms, integrated using a third-order Runge-Kutta/Matsuno time-stepping method.

\subsection{Reynolds Number}
The simulation starts by defining $Re$, and the physical parameters to conform to this, given a fixed background flow $U_{bg}$ are calculated. 
\begin{equation}
    Re = \frac{U_{bg} \cdot radius}{\nu}
\end{equation}
\textbf{Using $Re$, we solve for the required kinematic viscosity $\nu$ to satisfy the behavior dictated by $Re$.}
\begin{equation}
    \nu = \frac{U_{bg} \cdot radius}{Re}
\end{equation}

\subsection{Schmidt Number}
We also define $Sc=660$. Using the defined $Re$, we can solve for the Péclet number $Pe$, which is used more to verify our model is in the right physical realm. $Sc$ is the ratio of the kinematic viscosity to mass diffusivity:
\begin{equation}
    Sc = \frac{\nu}{K}
\end{equation}

Therefore, using the derived $\nu$ from above, and the predefined $Sc$:
\begin{equation}
    K_{diffusion} = \frac{\nu}{Sc}
\end{equation}
\textbf{Klawonn et al. used $K=0.0016 \space cm^2$}. In the ocean, dissolved gases such as $O_2$ diffuse orders of magnitude slower than momentum dissipates, so \textbf{$Sc \approx 660$}. 

\subsection{Péclet Number}
This is the ratio of advective to diffusive transport. 
\begin{equation}
Pe=Sc \cdot Re    
\end{equation}

\textbf{Kapellos et al. (2020)}: $0<Pe<1000$. As I said above, $Pe$ is used more to verify our model is in the right physical realm.

\subsection{Sherwood Number}
$Sh$ This is ratio of the total mass flux from the simulated particle to the purely diffusive flux in a stagnant medium. Essentially tells us the factor by which the sinking motion increases the nutrient or gas exchange compared to sitting still in that specific 2D setup.
\begin{equation}
Sh = 1 + 0.619Re \cdot 0.412 Sc ^{1/3}    
\end{equation}


In 2D, no standard value with advection turned off. For a \textbf{sphere} in a perfectly still fluid ($Re=0$), the Sherwood number is exactly 2, representing the absolute minimum mass transfer possible (pure spherical diffusion).


\subsection{Damköhler Number}
This is the ratio of the biological reaction rate to the \textbf{diffusion} rate.
\begin{equation}
    Da = \frac{k_{rem}/k_{Den1}/k_{Den2}...}{K}
\end{equation}
\textbf{Kapellos et al. (2020)}: $0<Da<10$. For example, for O2: If it is $\gg 1$, an anoxic core will form. If it is $< 1$, the particle will stay oxygenated and nothing will happen.




\subsection{Summary of Parameter Tradeoffs}
\begin{itemize}
    \item \textbf{Particle Radius ($r$)}: Defines the physical size of the biological aggregate.
        \begin{itemize}
            \item \textbf{Increase:} Increases the physical Reynolds number ($Re$), making the fluid wake prone to turbulent shedding. It also expands the anoxic core if respiration and diffusion rates remain constant.
            \item \textbf{Decrease:} Lowers the Reynolds number, yielding a highly stable, creeping flow wake. It increases oxygen penetration, potentially eliminating the anoxic core.
        \end{itemize}

    \item \textbf{Background Velocity ($U_{bg}$)}: The speed of the incoming water.
        \begin{itemize}
            \item \textbf{Increase:} Increases $Re$, triggering flow separation and von Kármán vortex shedding. \textit{Trade-off:} Requires a smaller advective time step (\texttt{dt$_{adv}$}) for stability, slowing computation.
            \item \textbf{Decrease:} Results in a stable, Stokes-like viscous flow without vortex shedding. Allows for larger advective time steps and faster computation.
        \end{itemize}

    \item \textbf{Kinematic Viscosity ($\nu$)}: The "stickiness" or momentum diffusivity of the fluid.
        \begin{itemize}
            \item \textbf{Increase:} Lowers $Re$ and suppresses vortex shedding, creating a wide, stable wake. This increases overall numerical stability.
            \item \textbf{Decrease:} Enables high-frequency vortex shedding due to sharper velocity gradients. \textit{Trade-off:} Demands higher grid resolution; otherwise, severe spatial gradients will cause the Arakawa solver to blow up.
        \end{itemize}

    \item \textbf{Mass Diffusivity ($K$)}: The rate at which tracer gases spread through the water.
        \begin{itemize}
            \item \textbf{Increase:} Artificially lowers the Schmidt number ($Sc$), blurring downstream chemical plumes and destroying tight vortex spirals. \textit{Trade-off:} Reduces the diffusive time step limit (\texttt{dt$_{diff}$}).
            \item \textbf{Decrease:} Preserves incredibly sharp, realistic chemical gradients and spirals within the wake vortices. \textit{Trade-off:} Slows oxygen penetration, expanding the anoxic core.
        \end{itemize}

    % \item \textbf{Aerobic Respiration Rate ($R$)}: The biological speed limit for $O_2$ consumption.
    %     \begin{itemize}
    %         \item \textbf{Increase:} Creates a high Damköhler number, resulting in a massive anoxic core due to rapid $O_2$ depletion before inward diffusion occurs. 
    %         \item \textbf{Decrease:} Shrinks the anoxic core by allowing deeper oxygen penetration. Has absolutely no effect on the physical fluid wake or hydrodynamics.
    %     \end{itemize}

    % \item \textbf{Denitrification Rate ($R_{denit}$)}: The biological production rate of $N_2O$ within the anoxic zone.
    %     \begin{itemize}
    %         \item \textbf{Increase:} Generates higher $N_2O$ concentrations, forming a darker, highly concentrated downstream plume.
    %         \item \textbf{Decrease:} Produces a fainter plume and must be precisely balanced against $R$ to maintain empirical $O_2$ to $N_2O$ emission ratios.
    %     \end{itemize}

    \item \textbf{Drag Penalty / Porosity ($\alpha$ / \texttt{drag\_mask})}: The force applied to stop fluid inside the solid boundary.
        \begin{itemize}
            \item \textbf{Increase:} Forces total flow separation, generating the shear stress required for vortex shedding. \textit{Trade-off:} The massive penalty term creates a highly stiff equation, severely restricting the drag time step limit (\texttt{dt$_{drag}$}).
            \item \textbf{Decrease:} Simulates a porous aggregate where water advects directly through the center, flushing out the anoxic core. It destroys flow separation, thereby killing vortex shedding.
        \end{itemize}
\end{itemize}

\section{Biological Processes}
\subsection{Oxygen}
To account for the biological consumption of oxygen within the particle, a respiration sink term, $R(x,y)$, is introduced. This term represents constant consumption and operates only within the particle. Key assumption: the particle has essentially unlimited organic matter that denitrifying bacteria can consume. The complete advection-diffusion-reaction equation for the tracer therefore becomes:
\begin{equation}
    \frac{\partial [O_2]}{\partial t} = -\vec{v} \cdot \nabla [O_2] + K \nabla^2 [O_2] - R_{resp}
\end{equation}
As the simulation approaches a steady state ($\frac{\partial C}{\partial t} \approx 0$), the profile of oxygen within the particle is determined strictly by the balance between the diffusive flux of oxygen penetrating the particle boundaries and the internal biological consumption rate $R(x,y)$.

\begin{itemize}
    \item \textbf{The size of the anoxic core was dependent on the O$_2$ concentration in the ambient water, that is, the anoxic interior expanded from 5\% of the total aggregate volume as an aggregate was suspended in 100\% air-saturated water to 95\% in 30\% air-saturated water.}
    \item There's 2 ways to change how much the tracer (O$_2$) actually penentrates into the particle. Increase K (which affects how well the tracers are CARRIED by the water), and decreasing R. Note: The radius of the particle also affects the time it takes for diffusion of tracers into the particle. Decreasing the drag mask also affects time it takes for tracers to penentrate, but it should be high enough that it acts like a solid.
\end{itemize}


\subsection{N_2\text{O}}

Similarly, the transport equation for nitrous oxide ($\text{N}_2\text{O}$) follows the exact same advection-diffusion physics but acts in reverse biologically. Instead of a sink, it has a denitrification source term ($S_{denit}$) that produces gas:
\begin{equation}
    \frac{\partial [\text{N}_2\text{O}]}{\partial t} = -\vec{v} \cdot \nabla [\text{N}_2\text{O}] + K \nabla^2 [\text{N}_2\text{O}] + S_{denit}
\end{equation}
Where the denitrification source term is modeled as a linear ramp activated below a critical oxygen threshold ($C_{thresh}$):
\begin{equation}
    S_{denit} = R_{denit} \cdot \max\left(0, \frac{[O_2]_{thresh} - [O_2]}{[O_2]_{thresh}}\right)
\end{equation}


Key assumption: the particle has essentially unlimited organic matter that denitrifying bacteria can consume.
\textbf{Why is there N2O in spots where oxygen is above the threshold. this is because it is DIFFUSING outward. Contributes to GHGs.}

\begin{itemize}
    \item \textbf{The ratio of O$_2$ consumption to N$_2$O production was 480:1.}


\end{itemize}




\section{CHANGE: Integrating NitrOMZ/MicrOMZ into Particle Model}
Instead of using the above predefined sms (source minus sink) equations, what if I use the sms from NitrOMZ/MircrOMZ? What would this entail?:
\begin{itemize}
    \item Obviously, oxygen levels are dictated by advection/diffusion of the flowing water. This informs RemOx.
    \item \textbf{Flow brings both O$_2$ and NO$_3$.} There is DOC ONLY within the particle.
    \item If I turn off advection, it turns into a 1D (rather than 2d) model (can look simply at one radius cross section)
\end{itemize}

\section{Comparison with Klawonn et al. (2015)}
\subsection{Klawonn's Results}
First, what results \textbf{could} I try to replicate?
\begin{itemize}
    \item Oxygen \& Anoxia Thresholds
    \begin{itemize}
        \item The anoxic core took up only about \textbf{5\% of the aggregate's total volume} in fully oxygenated water. In low-oxygen water (30\% air-saturation), the anoxic zone was \textbf{>95\% of the volume}.
        \item In fully aerated surface waters (250 $\mu$M $O_2$), cyanobacterial aggregates needed to be at least \textbf{1 mm or larger} to develop an anoxic center. In low oxygen waters (25 $\mu$M $O_2$), smaller aggregates (between \textbf{0.1 mm and 1 mm}) could turn anoxic.
    \end{itemize}
    
    \item $N_2O$ \& Denitrification
    \begin{itemize}
        \item The ratio of O$_2$ consumed to N$_2$O released at the aggregate surface was exactly \textbf{480:1}.
        \item When the flow of water around the aggregate was stopped, $N_2O$ accumulated fast in the center. Peak internal concentration: \textbf{0.9 to 2.2 $\mu$M}.
        \item Net $N_2$ production via denitrification was about \textbf{100 times lower} than the rate of nitrate reduction to ammonium.
    \end{itemize}
    
    \item $NH_4^+$ \& Nitrogen Fixation
    \begin{itemize}
        \item Most of the ammonium released by the aggregates (\textbf{88\%}) was driven by ammonification. Only \textbf{12\%} was directly from $N_2$ fixation.
        \item The cyanobacteria leaked a massive amount of the new nitrogen they fixed. \textbf{34.8 $\pm$ 21.8\%} of their $N_2$ fixation was released into the surrounding water as ammonium.
        \item Concentration of $NH_4^+$ was high, up to \textbf{35 $\mu$M in the center}, compared to a \textbf{0.52 $\mu$M in the ambient water}.
    \end{itemize}
\end{itemize}

\textbf{Some results from my simulation so far}:
\begin{itemize}
    \item Note: Increasing DOC doesn't increase relative consumption time. This is because of the kinetics that's a function of concentration.
    \item Most of the N compounds look ``normal'', as in when O$_2$ is low, NO$_3$ is consumed, and it goes down the denitrification chain. HOWEVER, we see an interesting ring looking at the N$_2$O plot. N$_2$O is actually being converted to N$_2$ quite fast, but that process is \textbf{very} inhibited by oxygen, so it doesn't happen in the out rim where oxygen is diffusing in. This gives us the cool looking ring.
    \item Right now, my viscosity is quite low and my diffusion is a bit high.
\end{itemize}


\section{ToDo}
\begin{itemize}
    \item What scientific questions can I ask? (see above)
    \item I think I'll need different diffusion constants for each gas, technically.
    \item Better parameterize oxygen/gas saturation? So right now, I just have an initial value. The units of these initial values in the code are in $mmol/m^3$. \textbf{I have to define a temperature to make the value as representative and physically meaningful as possible. Klawonn et al. (2015) measures oxygen concentration as \% saturated}
    \item Match the ambient water concentrations to be more realistic. For example, there's ammonium in water (e.g., 0.52 $\mu mol/L$).
    \item IMPORTANT!!: We technically need another process that converts POC (the particle) to DOC (the starting tracer we're using at the moment). \textbf{Can it be parameterized as a DOC flux instead?}
    \item MicrOMZ uses more first-principles (like thermodynamics). It will inevitably be better than NitrOMZ.
    \item Kapellos et al. (2022) is extremely helpful. They are also modeling slow-moving sinking particles and how bacteria affect their dissolution. I believe I can get inspiration from of lot of their physics.
\end{itemize}


\section{References}
\begin{enumerate}
    \item Nicholas K.-R Kevlahan, Jean-Michel Ghidaglia, Computation of turbulent flow past an array of cylinders using a spectral method with Brinkman penalization, \textit{European Journal of Mechanics - B/Fluids}, Volume 20, Issue 3, 2001, Pages 333-350.
    \item Klawonn, I., Bonaglia, S., Brüchert, V. et al. Aerobic and anaerobic nitrogen transformation processes in N2-fixing cyanobacterial aggregates. \textit{ISME J} 9, 1456–1466 (2015). https://doi.org/10.1038/ismej.2014.232
    \item Wan, X. S., et al. (2023): "Particle-associated incomplete denitrification is the primary source of N2O emissions in oxic waters" (Nature Communications). 
    \item Kapellos, G. E., Eberl, H. J., Kalogerakis, N., Doyle, P. S., \& Paraskeva, C. A. (2022). Impact of Microbial Uptake on the Nutrient Plume around Marine Organic Particles: High-Resolution Numerical Analysis. Microorganisms, 10(10), 2020. https://doi.org/10.3390/microorganisms10102020
\end{enumerate}




\end{document}
